\chapter{Einleitung}
\label{Einleitung}
\section{Motivation}

Wie kommt man überhaupt auf die Idee diese Thematik zu untersuchen?

\section{Problem Statement}

\section{Stand der Forschung}

Natürlich habe ich das Rad nicht erfunden, sondern ergänze nur die
Schraube der 783sten Speiche.  Daher muss ich aber auch die anderen
Speichen und besonders die Rad-Erfinder und die Geschichte der
Erfindung aufschreiben. Das ganze wird im Kapitel
\ref{StandDerForschung} genau betrachtet, hier wird es vorgreifend
zusammengefasst.

\section{Ansatz}

\section{Aufbau der Arbeit}
Im folgenden Kapitel \ref{StandDerForschung} werden Grundlagen
dargelegt und der Stand der Forschung diskutiert. In Kapitel
\ref{Ansatz} wird der Ansatz genauer beleuchtet. Im Kapitel
\ref{Implementierung} wird die Umsetzung des Ansatzes
beschrieben. Kapitel \ref{Validierung} praesentiert Resultate und eine
Diskussion der Ergebnisse. Im Kapitel \ref{ZusammenfassungAusblick}
wird die Arbeit zusammengefasst, ein Fazit gezogen und ein Ausblick
auf zukuenftige Arbeiten gegeben.


